% 第 6 课：板书第 1--3 张（后续照片继续追加于文件末尾）。
% 按老师推导顺序整理；对可确认的板书笔误作了统一修正，见文末「板书勘误」。
\section*{第 6 课}

% 第 1 张板书：积分定义、缩放及上半圆围道。
\paragraph{例 1.}
求实积分
\[
  I_{r,n}=\int_0^{+\infty}\frac{dx}{x^{2n}+r^{2n}},
  \qquad r>0,\quad n\in\mathbb{N}^{+}.
\]

\paragraph{解.}
作变量代换 $t=x/r$，得
\begin{align*}
  I_{r,n}
  &=\frac{1}{r^{2n}}\int_0^{+\infty}
    \frac{dx}{1+(x/r)^{2n}}\\
  &=\frac{1}{r^{2n-1}}\int_0^{+\infty}
    \frac{dt}{1+t^{2n}}
   =\frac{1}{r^{2n-1}}I_{1,n}.
\end{align*}
因此只需求 $I_{1,n}$。

令 $R>1$，取上半平面内的正向闭围道
\[
  \Gamma_R=[-R,R]\cup C_R,\qquad
  C_R:\ z=Re^{i\varphi},\quad 0\le\varphi\le\pi.
\]
由留数定理，
\begin{align*}
  \oint_{\Gamma_R}\frac{dz}{1+z^{2n}}
  &=\int_{-R}^{R}\frac{dx}{1+x^{2n}}
    +\int_{C_R}\frac{dz}{1+z^{2n}}\\
  &=2\pi i\sum_{k=1}^{n}
    \operatorname{Res}\left(\frac{1}{1+z^{2n}},z_k\right),
  \tag{*}
\end{align*}
其中 $z_1,\ldots,z_n$ 为上半平面内的极点。

% 第 2 张板书：大圆弧消失、极点位置及留数求和。
分子次数为 $0$，分母次数为 $2n\ge2$（$n\ge1$），因此当
$R\to+\infty$ 时，上半圆弧上的积分趋于 $0$。具体地，
\[
  \left|\int_{C_R}\frac{dz}{1+z^{2n}}\right|
  \le\frac{\pi R}{R^{2n}-1}\xrightarrow[R\to+\infty]{}0.
\]

在式 $(*)$ 中，极点满足
\[
  z^{2n}=-1=e^{(2k-1)\pi i}.
\]
取上半平面内的 $n$ 个根：
\[
  z_k=e^{\frac{(2k-1)\pi i}{2n}},
  \qquad k=1,2,\ldots,n.
\]
它们均为一阶极点，所以
\[
  \operatorname{Res}\left(\frac{1}{1+z^{2n}},z_k\right)
  =\frac{1}{2n z_k^{2n-1}}.
\]
由于 $z_k^{2n}=-1$，有
\[
  \frac{1}{z_k^{2n-1}}=-z_k.
\]
令 $R\to+\infty$，从式 $(*)$ 得
\begin{align*}
  2I_{1,n}
  &=2\pi i\sum_{k=1}^{n}\frac{1}{2n z_k^{2n-1}}\\
  &=\frac{\pi i}{n}\sum_{k=1}^{n}\frac{1}{z_k^{2n-1}}
   =-\frac{\pi i}{n}\sum_{k=1}^{n}z_k.
\end{align*}

令
\[
  \theta=\frac{\pi}{2n},\qquad
  q=e^{2i\theta}=e^{\pi i/n},
\]
则
\[
  z_k=\frac{q^k}{e^{i\theta}},
  \qquad q^n=e^{\pi i}=-1.
\]

% 第 3 张板书：等比数列求和及欧拉公式化简。
因此
\begin{align*}
  I_{1,n}
  &=-\frac{\pi i}{2n e^{i\theta}}\sum_{k=1}^{n}q^k\\
  &=-\frac{\pi i}{2n e^{i\theta}}\,
    \frac{q(1-q^n)}{1-q}\\
  &=-\frac{\pi i}{n}\,
    \frac{e^{i\theta}}{1-e^{2i\theta}}\\
  &=\frac{\pi i}{n}\,
    \frac{e^{i\theta}}{e^{2i\theta}-1}\\
  &=\frac{\pi}{2n\sin\theta},
\end{align*}
其中使用欧拉公式
\[
  \sin\theta=\frac{e^{i\theta}-e^{-i\theta}}{2i}
            =\frac{e^{2i\theta}-1}{2i e^{i\theta}}.
\]
最终
\[
  \boxed{
    I_{r,n}=\frac{\pi}
    {2n\,r^{2n-1}\sin\bigl(\frac{\pi}{2n}\bigr)}
  },\qquad r>0,\quad n\in\mathbb{N}^{+}.
\]

\paragraph{板书勘误（第 2--3 张）}
第 2 张的 $q^n$ 处写作 $+1$，应为 $-1$，因为
$q=e^{\pi i/n}$，从而 $q^n=e^{\pi i}=-1$。
另两张板书中的 $\theta$ 取值前后不一致；上文统一取
$\theta=\pi/(2n)$，并相应写为 $q=e^{2i\theta}$，
避免与 $q=e^{\pi i/n}$ 冲突。

% 此后继续追加第 6 课的后续板书，不另起导言区。


% 第 4 张板书：续前一例，用二倍角公式完成三角恒等式化简。
\paragraph{例 1 的另一种化简（续）}
仍令 $\theta=\pi/(2n)$。由上一页的等比数列求和结果，
\begin{align*}
  I_{1,n}
  &=-\frac{\pi i}{n}\frac{e^{i\theta}}{1-e^{2i\theta}}\\
  &=-\frac{\pi i}{n}
    \frac{e^{i\theta}}{1-\cos2\theta-i\sin2\theta}\\
  &=-\frac{\pi i}{n}
    \frac{e^{i\theta}}{2\sin^2\theta-2i\sin\theta\cos\theta}\\
  &=-\frac{\pi i}{2n}
    \frac{e^{i\theta}}{\sin\theta(\sin\theta-i\cos\theta)}\\
  &=-\frac{\pi i}{2n}
    \frac{e^{i\theta}}{(-i)(\cos\theta+i\sin\theta)\sin\theta}\\
  &=\frac{\pi}{2n}\frac{e^{i\theta}}{\sin\theta\,e^{i\theta}}
    =\frac{\pi}{2n\sin\theta}
    =\frac{\theta}{\sin\theta}
    =\frac{\pi}{2n\sin\frac{\pi}{2n}}.
\end{align*}
因此
\[
  I_{r,n}=\frac{1}{r^{2n-1}}I_{1,n}
  =\frac{\pi}{2nr^{2n-1}\sin\frac{\pi}{2n}}.
\]

% 第 4 张板书下半部：新例题、变量代换。
\paragraph{例 2.}
求实积分
\[
  J_{r,n}=\int_0^{+\infty}\frac{dx}{(x^2+r^2)^n},
  \qquad r>0,\quad n\in\mathbb N^{+}.
\]
令 $t=x/r$，则
\begin{align*}
  J_{r,n}
  &=\frac{1}{r^{2n}}
    \int_0^{+\infty}\frac{dx}{\bigl(1+(x/r)^2\bigr)^n}\\
  &=\frac{1}{r^{2n-1}}
    \int_0^{+\infty}\frac{dt}{(1+t^2)^n}
   =\frac{1}{r^{2n-1}}J_{1,n}.
\end{align*}
只需求
\[
  J_{1,n}=\int_0^{+\infty}\frac{dx}{(1+x^2)^n}.
\]

% 第 5 张板书：上半圆围道与 n 阶极点 i 的留数。
取上半圆正向闭围道
\[
  \Gamma_R=[-R,R]\cup C_R,\qquad R>1,\qquad
  C_R:\ z=Re^{i\varphi},\quad 0\le\varphi\le\pi.
\]
由留数定理，
\begin{align*}
  \oint_{\Gamma_R}\frac{dz}{(1+z^2)^n}
  &=\int_{-R}^{R}\frac{dx}{(1+x^2)^n}
    +\int_{C_R}\frac{dz}{(1+z^2)^n}\\
  &=2\pi i\,\operatorname{Res}
    \left(\frac{1}{(1+z^2)^n},i\right).
\end{align*}
当 $R\to+\infty$ 时，
\[
  \int_{C_R}\frac{dz}{(1+z^2)^n}\longrightarrow0.
\]
于是
\[
  \int_{-\infty}^{+\infty}\frac{dx}{(1+x^2)^n}
  =2\pi i\,\operatorname{Res}
   \left(\frac{1}{(1+z^2)^n},i\right)
  =2J_{1,n},
\]
即
\[
  J_{1,n}=\pi i\,\operatorname{Res}
  \left(\frac{1}{(1+z^2)^n},i\right).
\]

$\dfrac{1}{(1+z^2)^n}$ 在 $z=i$ 处有 $n$ 阶极点，
故要求其 Laurent 展开式中 $(z-i)^{-1}$ 的系数 $C_{-1}$。
将原函数写成
\begin{align*}
  \frac{1}{(1+z^2)^n}
  &=\frac{1}{(z-i)^n(z+i)^n}\\
  &=\frac{1}{(z-i)^n}
    \frac{1}{\bigl(z-(i-2i)\bigr)^n}\\
  &=\frac{1}{(z-i)^n}\cdot\frac{1}{(2i)^n}
    \cdot\frac{1}{\left(1+\dfrac{z-i}{2i}\right)^n}.
\end{align*}
所以只需求 $(1+X)^{-n}$ 展开式中 $X^{n-1}$ 的系数，
其中 $X=(z-i)/(2i)$。

% 第 6 张板书：广义二项式展开与负整数次幂系数。
\paragraph{引理（广义二项式展开）}
当 $|X|<1$ 时，
\[
  (1+X)^\alpha
  =1+\sum_{k=1}^{+\infty}
    \frac{\alpha(\alpha-1)\cdots(\alpha-k+1)}{k!}\,X^k.
\]
令 $\alpha=-n$，其中 $n\in\mathbb N^{+}$，得
\begin{align*}
  \frac{1}{(1+X)^n}
  &=1+\sum_{k=1}^{+\infty}
    \frac{(-n)(-n-1)\cdots(-n-k+1)}{k!}\,X^k\\
  &=1+\sum_{k=1}^{+\infty}
    (-1)^k\frac{(n+k-1)!}{k!(n-1)!}\,X^k\\
  &=1+\sum_{k=1}^{+\infty}
    (-1)^k\binom{n+k-1}{n-1}X^k\\
  &=\sum_{k=0}^{+\infty}
    (-1)^k\binom{n+k-1}{n-1}X^k.
\end{align*}
% 第 6 张板书到此为止。后续照片继续求 C_{-1}，不提前添加解答。
